\section{Non-Linear Deliberation Workflow}
\label{sec:pa-deliberation}

\subsection{Automatic Rework}
\label{sec:pa-rework}
When the analyst's result is not OK, because issues need a stakeholder's answer, a control node writes a \emph{rework directive}: a record that names the open ambiguity issues and conflicts to settle and the personas who can settle them. The interview continues with these personas; the interviewer reads the directive first and asks about the issues, and the new messages are appended to the same transcripts. The analyst then runs again on the current analysis and resolves the issues the new answers settle. The number of rework loops is bounded by the maximum review rounds.

\subsection{Human Review Checkpoints}
\label{sec:pa-checkpoints}
At a checkpoint the run pauses and waits for the user (Table~\ref{tab:pa-checkpoints}).

\begin{table}[H]
    \centering
    \caption{Human review checkpoints and their options.}
    \label{tab:pa-checkpoints}
    \begin{tabularx}{\linewidth}{|l|X|X|}
        \hline
        \textbf{Checkpoint} & \textbf{Reviewed artifact} & \textbf{Options and next node} \\
        \hline
        Finalize the interviewees & Persona catalog, editable & Approve (after edits) $\rightarrow$ interviewer \\
        Review the analysis & Requirements analysis, editable & Accept $\rightarrow$ requirement engineer; reject and re-analyse $\rightarrow$ analyst; reject and re-interview $\rightarrow$ rework elicitation \\
        Review the specification & Requirement specification, editable & OK $\rightarrow$ backlog manager; not OK $\rightarrow$ requirement engineer \\
        \hline
    \end{tabularx}
\end{table}

Every review also offers an option \emph{Something else}: it asks the user for words and sends them to the intent router as a message, so a decision that no option expresses is still possible.

\subsection{Incremental Intent Routing}
\label{sec:pa-intent-router}
A message the user sends while the run works is taken in at the next safe point, after a node finishes: the intent router decides which node should handle it (planner, interviewer, analyst, requirement engineer or backlog manager) and the run continues there. Every run starts at the intent router, the first one included; a project that holds nothing yet is sent on to the planner.

The intent router reads the pending messages, newest first, together with the current specification and backlog, chooses the node whose artifact the message changes, and restates the intent as one instruction to that node. Some examples:
\begin{itemize}
    \item ``Also support booking from a smartwatch'' adds a capability that no stakeholder discussed: it is routed to the \emph{interviewer}, so that the personas are asked about it.
    \item ``Students must not book more than two rooms per week'' is a decision of the product owner rather than a stakeholder need: it is routed to the \emph{requirement engineer}, which adds a constraint.
    \item ``Split the booking story into booking and cancelling'' concerns only the backlog: it is routed to the \emph{backlog manager}.
\end{itemize}
A message that only restates a wish already covered by the artifacts is routed to the node that would run next. The router never routes to a persona agent, since the user does not speak in the place of a stakeholder.

\subsection{Bounded Automatic Decisions}
\label{sec:pa-decisions}

Three gates of the workflow take a decision that is not a free generation: whether the interview of a persona should stop, whether a specified requirement has the right type and priority, and whether validation passed. At these gates the decision can be taken by a separate, calibrated decision model instead of the agent alone. When the confidence of the decision model is high, its decision is applied automatically; between two thresholds the agent's own value is kept; below the low threshold the run asks the user through a review. The decision model never changes the workflow itself: it only decides which of the existing paths is taken. The decision layer is off by default and is used only when a decision provider is configured; whether it is used in the experiments is stated with the experimental setup (Section~\ref{sec:eval-details}).

\subsection{Questions About the Application}
\label{sec:pa-assistance}
Not every message of the user is about the requirements. A user may ask where the interview with a persona is, what the cost limit of the project is, or why the last run failed. Such a message should not start a run or reach the intent router. For this reason every message is first given to an \emph{assistant}, the single agent of a second, separate workflow, which decides whether the message belongs to the requirements work or to the application:
\begin{itemize}
    \item A message about the work enters the requirements engineering workflow exactly as described above: it starts a run, or waits as pending in the run that is unfinished.
    \item A message about the application is answered by the assistant. It sees a summary of the project and may read what it is missing through eleven read-only tools: an overview of the project, the list of its documents, a document, a search over the items, an item, the evidence of an item, the open findings, the last run, the run settings, a step of the workflow and the budget of the month. The number of reads per message is bounded (five by default). The answer is a few sentences, stored as a message of the thread, and may carry up to four kinds of command for the web app: open a document, open an item, go to the pending review, or open the run settings.
\end{itemize}
The assistant never writes an artifact, never starts a run of the requirements workflow and never changes its graph; its own answer is recorded as a short run of the assistance workflow, so that the web app shows it like any other work of an agent. Every command it names is checked against a catalog of commands before it reaches the web app, and when the assistant fails or is not certain, the message goes to the requirements workflow, so that no message is lost. The cost of each answer is recorded and counts against the budget of the user like the cost of a run.
